\scanpage{87}
% archive sheet number: 86
On peut poser à l'autre extrême

\textcircled{$d''$} $\varprojlim$ de théories (sans conditions particulières) \ldots\ Il suffirait de le postuler en un cas d'indices bien ordonné\unsure{s|s.|e}\ldots

\uline{Remarque 1}\quad On ne sort à aucun moment des théories $T$ 2 représentables par $B_T$, où $B_T=R$ \deleted{\unsure{si} \unsure{eng.}}\inserted{admet une petite} \deleted{famille} \unsure{gén.}, $I\to\mathrm{Ob}\,R$, i.e. telle que $\forall C\in\mathrm{Ob}\,\Delta$, $\underline{\mathrm{Hom}}_{\Delta}(R,C)\xrightarrow[\mathrm{fidèle}]{}\deleted{\underline{\mathrm{Hom}}}C^I$

\deleted{Remarque}

\uline{Remarque 2}\quad Soit réciproquement $R\in\mathrm{Ob}\,\Delta$, avec $\Delta$-gén. $I\to\mathrm{Ob}\,R$. \deleted{On va procéder \unsure{Soit Munissons}} \deleted{$T_0$ repr par $\unsure{\tau} R_0=\widetilde I$}

$1^\circ)$ $R$ \uline{petite}. On prend $I=\mathrm{Ob}\,R$, \unsure{qui|on|ou} part de
\[
T_0=\tau^I\ ;\ B'_{T_0}=\widetilde I=R_0
\]
On prend ensuite $H=\mathrm{Fl}(R)$, qu'on envoie dans $\mathrm{Ob}\,R_0\times\mathrm{Ob}\,R_0$ par le composé
\[
 H\longrightarrow\underset{\substack{\Vert\\I\times I}}{\mathrm{Ob}\,R\times\mathrm{Ob}\,R}
 \longrightarrow\mathrm{Ob}\,R_0\times\mathrm{Ob}\,R_0
\]
d'où $\underset{\unsure{A}}{T_1},R_1$ avec $\underset{\substack{\uparrow\\ I}}{R_0}\mathbin{\unsure{-|\to}}R_1$, $T_1\xrightarrow[\deleted{\mathrm{fid}}]{}T_0$.

% slanted; brace alongside the first three stages.
\margin{$\Biggl\{$ les trois premiers pas montrent que si $J$ est petite, la théorie $C\mapsto\underline{\mathrm{Hom}}_{\mathrm{Cat}}(J,C)$ (2-représenté par $\widetilde J$) est une théorie alg.}

On exprime ensuite que la \inserted{composition des} flèches de $R$ devient la comp. des flèches de $R_1$ i.e. \deleted{l'application} ``préfoncteur'' $R\to R_1$ devient un foncteur. On trouve ainsi
\[
T_2,R_2,\qquad R_1\mathbin{-}R_2,\quad\text{où }R_2=\widetilde R
\quad(T_2(C)\simeq\mathrm{Hom}_{\mathrm{Cat}}(R,C))
\]
On exprime enfin que le foncteur $R\to R_2$ devient un $\Delta$-foncteur, \deleted{\unsure{en}} [i.e commute à certains types de lim] en écrivant que certaines flèches de $R_2$ deviennent inversibles, d'où
\[
 T_3,R_3\qquad R_3=R .
\]
% slanted
\margin{On utilise ici que les $\Delta$-foncteurs $=$ à ceux qui commutent à certains types de lim\ldots\ Sinon il faut exprimer une condition}

[\uline{NB} On a appliqué 1 fois pas a), une fois pas c), deux fois pas b), et pas du tout d)]

$2^\circ)$ $R$ quelconque, mais avec $I\mathbin{-}\mathrm{Ob}\,R$ ``générateurs''. Soit $\rho$ la sous catégor. pleine de $R$ \uline{engendrée} par $I$. On a
\[
\underline{\mathrm{Hom}}_{\Delta}(R,C)\xrightarrow{\mathrm{fid}}
\underbrace{\underline{\mathrm{Hom}}_{\mathrm{Cat}}(\rho_0,C)}_{\text{théorie alg.}}
 \simeq\underline{\mathrm{Hom}}_{\unsure{\Delta|\mathrm{Cat}}}(\widetilde\rho_0,R)
\]
% slanted; third numbered case climbs up from the lower left.
\margin{$3^\circ)$ Soit $J\xrightarrow{u}I$ une flèche \unsure{$\underline{\mathrm{Cat}}$} dans $(\mathrm{Cat})$, avec $J\in\mathrm{Ob}\,\Delta$. Alors $C\mapsto$ sous-catégorie de $\underline{\mathrm{Hom}}(I,C)$ formée des $f:I\to C$ tels que $f\circ u$ soit un $\Delta$ foncteur, est 2-représ.}
